\chapter{Étude préliminaire}

\section*{Introduction}
Ce chapitre pose les fondations analytiques du projet. Nous identifions les acteurs
du système, formalisons les besoins fonctionnels et non fonctionnels, construisons
le product backlog avec 28 user stories, puis décrivons l'architecture globale et
la pile technologique retenue.

% ============================================================
\section{Les acteurs du système}

\subsection{Rôles applicatifs}

\begin{description}[leftmargin=2.5cm]
  \item[\textbf{Admin}]   Accès complet : gestion des utilisateurs, équipes, audit,
        toutes les fonctionnalités. Supervise l'ensemble des projets.
  \item[\textbf{Member}]  Profil standard : gestion de ses projets, tâches, IA,
        Gantt, calendrier, notifications.
  \item[\textbf{Client}]  Accès en lecture seule aux projets auxquels il est invité.
        Peut consulter l'avancement et les livrables.
\end{description}

\subsection{Rôles dans un projet}

\begin{table}[H]
  \centering
  \caption{Rôles dans un projet et droits associés}
  \label{tab:roles_projet}
  \begin{tabularx}{\textwidth}{L{2.5cm}X}
    \toprule
    \thead{Rôle projet} & \thead{Droits} \\
    \midrule
    \rowcolor{lightgray}
    owner   & Toutes les actions sur le projet, y compris suppression et archivage \\
    manager & Modifier le projet, gérer les membres, créer/supprimer des tâches et statuts \\
    \rowcolor{lightgray}
    developer & Créer et modifier des tâches, ajouter des commentaires, saisir du temps \\
    viewer  & Lecture seule : consulter les tâches, le Gantt, les commentaires \\
    \bottomrule
  \end{tabularx}
\end{table}

\subsection{Acteurs logiciels}

\begin{description}[leftmargin=3.5cm]
  \item[\textbf{Service ML}] Microservice FastAPI (Python 3.11) exposant les endpoints
        de recommandation, prédiction de retard et estimation de durée.
  \item[\textbf{LLM OpenRouter}] API tierce (Llama 3.3 70B Instruct) invoquée pour
        la génération de justifications en langage naturel.
\end{description}

\subsection{Matrice des droits}

\begin{table}[H]
  \centering
  \caption{Matrice des droits par rôle système}
  \label{tab:matrice_droits}
  \small
  \begin{tabularx}{\textwidth}{L{5cm}C{2.5cm}C{2.5cm}C{2.5cm}}
    \toprule
    \thead{Action} & \thead{Admin} & \thead{Member} & \thead{Client} \\
    \midrule
    \rowcolor{lightgray}
    Gérer les utilisateurs          & \cmark & \xmark & \xmark \\
    Gérer les équipes               & \cmark & \xmark & \xmark \\
    \rowcolor{lightgray}
    Créer un projet                 & \cmark & \cmark & \xmark \\
    Accéder au Gantt                & \cmark & \cmark & \xmark \\
    \rowcolor{lightgray}
    Accéder au calendrier           & \cmark & \cmark & \xmark \\
    Accéder aux fonctionnalités IA  & \cmark & \cmark & \xmark \\
    \rowcolor{lightgray}
    Consulter ses projets           & \cmark & \cmark & \cmark \\
    Consulter le journal d'audit    & \cmark & \cmark & \xmark \\
    \rowcolor{lightgray}
    Recevoir des notifications      & \cmark & \cmark & \cmark \\
    Modifier son profil             & \cmark & \cmark & \cmark \\
    \bottomrule
  \end{tabularx}
\end{table}

% ============================================================
\section{Besoins fonctionnels}

\begin{description}[leftmargin=1.5cm]
  \item[\textbf{AUTH}] BF-AUTH-01 : s'inscrire et se connecter ; BF-AUTH-02 :
        réinitialiser le mot de passe par email ; BF-AUTH-03 : gérer son profil
        et ses préférences.
  \item[\textbf{PROJ}] BF-PROJ-01 : créer et configurer un projet ; BF-PROJ-02 :
        inviter des membres par email ; BF-PROJ-03 : gérer les statuts Kanban
        personnalisés.
  \item[\textbf{SPRINT}] BF-SPRINT-01 : créer et planifier un sprint ; BF-SPRINT-02 :
        démarrer et clôturer un sprint avec report automatique du backlog.
  \item[\textbf{TASK}] BF-TASK-01 : créer des tâches hiérarchiques ; BF-TASK-02 :
        définir des dépendances FS/SS/FF/SF ; BF-TASK-03 : affecter des membres
        avec allocation en pourcentage ; BF-TASK-04 : déplacer les tâches en
        Kanban par drag-and-drop.
  \item[\textbf{PLAN}] BF-PLAN-01 : visualiser le Gantt avec chemin critique ;
        BF-PLAN-02 : consulter le calendrier mensuel.
  \item[\textbf{NOTIF}] BF-NOTIF-01 : recevoir des notifications temps réel
        in-app ; BF-NOTIF-02 : configurer les préférences de notification.
  \item[\textbf{DASH}] BF-DASH-01 : consulter un dashboard personnalisé selon le
        rôle ; BF-DASH-02 : accéder au journal d'audit.
  \item[\textbf{IA}] BF-IA-01 : obtenir une recommandation ML d'assignation ;
        BF-IA-02 : estimer la durée d'une tâche par LLM ; BF-IA-03 : détecter
        les risques ; BF-IA-04 : utiliser le chat assistant IA.
\end{description}

% ============================================================
\section{Besoins non fonctionnels}

\begin{table}[H]
  \centering
  \caption{Besoins non fonctionnels}
  \label{tab:bnf}
  \begin{tabularx}{\textwidth}{L{3cm}L{5cm}X}
    \toprule
    \thead{Catégorie} & \thead{Exigence} & \thead{Valeur cible} \\
    \midrule
    \rowcolor{lightgray}
    Sécurité     & Auth stateless Bearer Token         & Sanctum ; expiration configurable \\
    Sécurité     & RBAC à 2 niveaux                   & Spatie Permission + ProjectPolicy \\
    \rowcolor{lightgray}
    Performance  & Latence P95 API                    & $\leq$ 500 ms \\
    Performance  & Inférence ML (P50)                 & $\leq$ 45 ms \\
    \rowcolor{lightgray}
    Ergonomie    & Responsive (breakpoints)           & Mobile, tablette, desktop \\
    Ergonomie    & Accessibilité                      & WCAG 2.1 niveau AA \\
    \rowcolor{lightgray}
    Fiabilité    & Disponibilité cible                & 99,5\,\% \\
    Fiabilité    & Détection cycle dépendances        & BFS avant toute insertion \\
    \rowcolor{lightgray}
    Maintenabilité & Duplication code (SonarCloud)    & $\leq$ 3\,\% \\
    Maintenabilité & Couverture tests (backend)       & $\geq$ 80\,\% \\
    \bottomrule
  \end{tabularx}
\end{table}

% ============================================================
\section{Product Backlog — User Stories}

\begin{longtable}{C{1.2cm}L{2cm}L{6.5cm}C{2cm}C{1.5cm}C{1.5cm}}
  \caption{Product Backlog — 28 User Stories}
  \label{tab:product_backlog}\\
  \toprule
  \thead{ID} & \thead{Acteur} & \thead{User Story} & \thead{MoSCoW} & \thead{Pts} & \thead{Sprint} \\
  \midrule
  \endfirsthead
  \multicolumn{6}{c}{\small\textit{(suite du tableau \ref{tab:product_backlog})}} \\
  \toprule
  \thead{ID} & \thead{Acteur} & \thead{User Story} & \thead{MoSCoW} & \thead{Pts} & \thead{Sprint} \\
  \midrule
  \endhead
  \bottomrule
  \endfoot

  \rowcolor{lightgray}
  US-01 & Tous    & S'inscrire et se connecter en sécurité              & Must    & 5  & S1 \\
  US-02 & Admin   & Gérer les comptes utilisateurs                      & Must    & 8  & S1 \\
  \rowcolor{lightgray}
  US-03 & Admin   & Définir rôles et permissions Spatie                  & Must    & 5  & S1 \\
  US-04 & Tous    & Modifier profil et préférences                      & Should  & 5  & S1 \\
  \rowcolor{lightgray}
  US-05 & Tous    & Réinitialiser mot de passe par email                & Should  & 3  & S1 \\
  US-06 & Member  & Créer et configurer un projet                       & Must    & 5  & S2 \\
  \rowcolor{lightgray}
  US-07 & Chef    & Inviter des membres par email (7 jours)              & Must    & 8  & S2 \\
  US-08 & Chef    & Définir des statuts Kanban personnalisés             & Should  & 5  & S2 \\
  \rowcolor{lightgray}
  US-09 & Chef    & Créer et démarrer un sprint                         & Must    & 8  & S2 \\
  US-10 & Chef    & Clôturer un sprint avec report du backlog           & Should  & 5  & S2 \\
  \rowcolor{lightgray}
  US-11 & Member  & Créer des tâches hiérarchiques (sous-tâches)        & Must    & 8  & S3 \\
  US-12 & Member  & Définir des dépendances FS/SS/FF/SF                 & Must    & 8  & S3 \\
  \rowcolor{lightgray}
  US-13 & Member  & Affecter des membres avec allocation \%             & Must    & 5  & S3 \\
  US-14 & Member  & Déplacer une tâche en Kanban (drag-and-drop)        & Must    & 5  & S3 \\
  \rowcolor{lightgray}
  US-15 & Member  & Saisir le temps passé sur une tâche                 & Should  & 3  & S3 \\
  US-16 & Member  & Commenter une tâche et mentionner (@)               & Should  & 5  & S3 \\
  \rowcolor{lightgray}
  US-17 & Member  & Visualiser le Gantt avec chemin critique            & Must    & 13 & S4 \\
  US-18 & Member  & Consulter le calendrier mensuel                     & Should  & 5  & S4 \\
  \rowcolor{lightgray}
  US-19 & Tous    & Recevoir des notifications temps réel               & Must    & 8  & S4 \\
  US-20 & Tous    & Configurer ses préférences de notification          & Should  & 3  & S4 \\
  \rowcolor{lightgray}
  US-21 & Member  & Accéder à un dashboard personnel                   & Should  & 8  & S4 \\
  US-22 & Admin   & Accéder au dashboard administrateur                 & Should  & 5  & S4 \\
  \rowcolor{lightgray}
  US-23 & Member  & Obtenir une recommandation ML d'assignation         & Must    & 13 & S5 \\
  US-24 & Member  & Estimer la durée d'une tâche par IA                 & Must    & 8  & S5 \\
  \rowcolor{lightgray}
  US-25 & Member  & Discuter avec l'assistant IA                        & Should  & 8  & S5 \\
  US-26 & Admin   & Relancer l'entraînement du modèle ML                & Could   & 5  & S5 \\
  \rowcolor{lightgray}
  US-27 & Client  & Consulter le rapport d'avancement de son projet     & Should  & 5  & S4 \\
  US-28 & Tous    & Utiliser la recherche globale                       & Should  & 5  & S4 \\
\end{longtable}

% ============================================================
\section{Diagramme global des cas d'utilisation}

\begin{figure}[H]
  \centering
  \IfFileExists{figures/ch2_global_use_case.png}{%
    \includegraphics[width=0.92\textwidth]{figures/ch2_global_use_case.png}%
  }{%
    \fbox{\parbox{0.92\textwidth}{\centering\vspace{2cm}
      \textbf{[IMAGE A INSÉRER]}\\[0.4em]
      \texttt{figures/ch2\_global\_use\_case.png}\\[0.3em]
      \textit{Diagramme global des cas d'utilisation — 3 acteurs humains, 2 acteurs logiciels,
      5 packages fonctionnels}
    \vspace{2cm}}}%
  }
  \caption{Diagramme global des cas d'utilisation — TAC-TIC IntelliPM}
  \label{fig:global_use_case}
\end{figure}

Le diagramme regroupe les 28 cas d'utilisation en cinq packages :
\textbf{Authentification}, \textbf{Projets \& Équipes}, \textbf{Tâches},
\textbf{Planification} et \textbf{Intelligence Artificielle}. Les acteurs logiciels
(Service ML FastAPI, LLM OpenRouter) interviennent en tant qu'acteurs secondaires
dans le package IA.

% ============================================================
\section{Conception architecturale}

\subsection{Architecture en couches}

\begin{figure}[H]
  \centering
  \IfFileExists{figures/ch2_architecture_couches.png}{%
    \includegraphics[width=0.85\textwidth]{figures/ch2_architecture_couches.png}%
  }{%
    \fbox{\parbox{0.85\textwidth}{\centering\vspace{1.5cm}
      \textbf{[IMAGE A INSÉRER]}\\[0.4em]
      \texttt{figures/ch2\_architecture\_couches.png}\\[0.3em]
      \textit{Architecture en 5 couches : Présentation, Services applicatifs, Métier,
      Persistance, IA}
    \vspace{1.5cm}}}%
  }
  \caption{Architecture en couches de TAC-TIC IntelliPM}
  \label{fig:archi_couches}
\end{figure}

\subsection{Architecture distribuée}

\begin{figure}[H]
  \centering
  \IfFileExists{figures/ch2_architecture_distribuee.png}{%
    \includegraphics[width=0.90\textwidth]{figures/ch2_architecture_distribuee.png}%
  }{%
    \fbox{\parbox{0.90\textwidth}{\centering\vspace{1.5cm}
      \textbf{[IMAGE A INSÉRER]}\\[0.4em]
      \texttt{figures/ch2\_architecture\_distribuee.png}\\[0.3em]
      \textit{Architecture distribuée : React SPA $\leftrightarrow$ Laravel API $\leftrightarrow$ MySQL + FastAPI ML}
    \vspace{1.5cm}}}%
  }
  \caption{Architecture distribuée — Vue déploiement}
  \label{fig:archi_distribuee}
\end{figure}

\subsection{Pile technologique}

\begin{table}[H]
  \centering
  \caption{Pile technologique complète}
  \label{tab:techstack}
  \small
  \begin{tabularx}{\textwidth}{L{3cm}L{3.5cm}C{1.8cm}X}
    \toprule
    \thead{Catégorie} & \thead{Outil} & \thead{Version} & \thead{Rôle} \\
    \midrule
    \rowcolor{lightgray}
    Frontend      & React                  & 19       & SPA, composants UI \\
    Frontend      & Vite                   & 7        & Bundler, HMR \\
    \rowcolor{lightgray}
    Frontend      & TailwindCSS            & 4        & Styles utilitaires \\
    Frontend      & TanStack React Query   & v5       & Cache et état serveur \\
    \rowcolor{lightgray}
    Frontend      & React Router DOM       & 6        & Routing SPA \\
    Frontend      & Framer Motion          & --       & Animations \\
    \rowcolor{lightgray}
    Backend       & Laravel                & 12       & Framework API REST \\
    Backend       & PHP                    & 8.2      & Langage serveur \\
    \rowcolor{lightgray}
    Backend       & Laravel Sanctum        & --       & Auth Bearer Token \\
    Backend       & Spatie Permission      & --       & RBAC système \\
    \rowcolor{lightgray}
    Backend       & Laravel Reverb         & --       & WebSocket server \\
    Backend       & Eloquent ORM           & --       & Mapping objet-relationnel \\
    \rowcolor{lightgray}
    Base de données & MySQL                & 8        & Persistance principale \\
    IA/ML         & FastAPI                & --       & Microservice ML \\
    \rowcolor{lightgray}
    IA/ML         & Scikit-learn           & 1.4      & GBM, RF, LogReg \\
    IA/ML         & OpenRouter (LLM)       & --       & Llama 3.3 70B Instruct \\
    \rowcolor{lightgray}
    Qualité       & SonarCloud             & --       & Analyse statique \\
    \bottomrule
  \end{tabularx}
\end{table}

\subsection{Diagramme de classes global}

\begin{figure}[H]
  \centering
  \IfFileExists{figures/ch2_class_diagram.png}{%
    \includegraphics[width=0.95\textwidth]{figures/ch2_class_diagram.png}%
  }{%
    \fbox{\parbox{0.95\textwidth}{\centering\vspace{2cm}
      \textbf{[IMAGE A INSÉRER]}\\[0.4em]
      \texttt{figures/ch2\_class\_diagram.png}\\[0.3em]
      \textit{Diagramme de classes global — 18 entités avec relations et cardinalités}
    \vspace{2cm}}}%
  }
  \caption{Diagramme de classes global de TAC-TIC IntelliPM}
  \label{fig:class_global}
\end{figure}

\begin{table}[H]
  \centering
  \caption{Les 18 entités du modèle de données}
  \label{tab:entites}
  \begin{tabularx}{\textwidth}{L{4cm}X}
    \toprule
    \thead{Classe} & \thead{Rôle} \\
    \midrule
    \rowcolor{lightgray}
    User            & Utilisateur du système (profil, auth, préférences) \\
    Team            & Équipe de développement (leader, membres) \\
    \rowcolor{lightgray}
    Specialty       & Compétence/technologie associée à un utilisateur \\
    Project         & Projet avec code, statut, dates, équipe \\
    \rowcolor{lightgray}
    ProjectInvitation & Invitation email (token UUID, expiration 7 jours) \\
    Sprint          & Sprint Scrum avec objectif, dates, statut \\
    \rowcolor{lightgray}
    TaskStatus      & Colonne Kanban (nom, couleur, ordre) \\
    Task            & Tâche avec hiérarchie, priorité, workflow, IA \\
    \rowcolor{lightgray}
    TaskAssignment  & Assignation avec pourcentage d'allocation \\
    TaskDependency  & Dépendance FS/SS/FF/SF entre tâches \\
    \rowcolor{lightgray}
    TaskComment     & Commentaire (threaded) avec mentions @ \\
    Attachment      & Pièce jointe (fichier ou lien) \\
    \rowcolor{lightgray}
    TimeEntry       & Temps passé (manuel ou timer) \\
    Notification    & Notification in-app avec type et données JSON \\
    \rowcolor{lightgray}
    ActivityLog     & Journal d'activité projet \\
    AuditLog        & Journal d'audit système (IP, user, action) \\
    \rowcolor{lightgray}
    AiRecommendation & Résultat d'une requête IA (ML ou LLM) \\
    PersonalAccessToken & Token Sanctum pour l'authentification \\
    \bottomrule
  \end{tabularx}
\end{table}

\subsection{Schéma ERD}

\begin{figure}[H]
  \centering
  \IfFileExists{figures/ch2_erd.png}{%
    \includegraphics[width=0.95\textwidth]{figures/ch2_erd.png}%
  }{%
    \fbox{\parbox{0.95\textwidth}{\centering\vspace{2cm}
      \textbf{[IMAGE A INSÉRER]}\\[0.4em]
      \texttt{figures/ch2\_erd.png}\\[0.3em]
      \textit{Schéma ERD — 42 tables avec clés primaires, étrangères et contraintes}
    \vspace{2cm}}}%
  }
  \caption{Schéma Entity-Relationship Diagram (ERD) — Base de données MySQL}
  \label{fig:erd}
\end{figure}

\section*{Conclusion}

L'étude préliminaire a permis de formaliser l'ensemble des exigences du système :
28 user stories priorisées en MoSCoW, réparties sur 5 sprints pour un total de
177 story points. L'architecture distribuée en trois couches (React~SPA,
Laravel~API, FastAPI~ML) garantit l'indépendance de déploiement et la scalabilité
de chaque composant. Les chapitres suivants détaillent la réalisation sprint par
sprint.
